% GENERATED FILE. Edit canonical lessons, then run npm run generate:books.
% canonical-source-hash: fnv1a64:91b76e12e4561312
% canonical-lessons: HI-C322-talvar, HI-C322-dhal, HI-C322-bhala, HI-C322-tir, HI-C322-prayah

\chapter{Swords and Shields}
\label{ch:hi-swords-and-shields}
\hlchaptermodality{322}

\begin{chapteropening}
\textbf{By the end of this chapter:} \emph{I can describe how battles were fought long ago.}

Complete the last lesson of chapter 322: I can describe how battles were fought long ago.
\end{chapteropening}

\section[\texorpdfstring{\hi{तलवार} (talvār)}{तलवार (talvār)}]{\hi{तलवार} (talvār) — a sword}

\label{lesson:HI-C322-talvar}



\begin{quote}
Before the new one: say the Hindi for to clash, then the Hindi for to fight hard.
\end{quote}

\subsection*{You'll want to know: \hi{तलवार}}
\textbf{\hi{तलवार}} — \emph{talvār} — \textquotedblleft{}a sword\textquotedblright{}.

\begin{quote}
\hi{राजा} \hi{ने} \hi{तलवार} \hi{उठाई।}
\end{quote}

\begin{quote}
\emph{rājā ne talvār uṭhāī.}
\end{quote}

\begin{quote}
The king raised his sword.
\end{quote}

Feminine.

\begin{cousinweb}
\textbf{\hi{तलवार}} (\emph{talvār}) goes back to Sanskrit \emph{taravāri}, \textquotedblleft{}a sword\textquotedblright{} (CDIAL 5706), which CDIAL gives with a Hindi \emph{talwār}. It is feminine: \textbf{\hi{तलवार} \hi{उठाई}}.
\end{cousinweb}

\begin{grammarlens}[title={What you've built}]
One more word for telling what happened.
\end{grammarlens}

\subsection*{Your turn}
\begin{itemize}
\raggedright
  \item \emph{Say it:} \emph{talvār}
  \item \emph{Say it:} \emph{talvār}, once more
  \item \emph{Say it:} \emph{rājā ne talvār uṭhāī.} — The king raised his sword.
  \item \emph{Recall:} say the Hindi for to clash, then the Hindi for to fight hard, then say \emph{talvār} again
\end{itemize}

\begin{tcolorbox}[breakable,colback=teal!4,colframe=teal!35!black,title={Before you move on}]
What does \hi{तलवार} mean? (\textquotedblleft{}A sword\textquotedblright{}.) Say it once more.
\end{tcolorbox}
\section[\texorpdfstring{\hi{ढाल} (ḍhāl)}{ढाल (ḍhāl)}]{\hi{ढाल} (ḍhāl) — a shield}

\label{lesson:HI-C322-dhal}



\begin{quote}
Before the new one: say the Hindi for to fight hard, then the Hindi for a sword.
\end{quote}

\subsection*{You'll want to know: \hi{ढाल}}
\textbf{\hi{ढाल}} — \emph{ḍhāl} — \textquotedblleft{}a shield\textquotedblright{}.

\begin{quote}
\hi{सिपाही} \hi{के} \hi{एक} \hi{हाथ} \hi{में} \hi{तलवार} \hi{और} \hi{दूसरे} \hi{में} \hi{ढाल} \hi{थी।}
\end{quote}

\begin{quote}
\emph{sipāhī ke ek hāth meṁ talvār aur dūsre meṁ ḍhāl thī.}
\end{quote}

\begin{quote}
The soldier had a sword in one hand and a shield in the other.
\end{quote}

Feminine.

\begin{cousinweb}
\textbf{\hi{ढाल}} (\emph{ḍhāl}) goes back to Sanskrit \emph{ḍhāla}, \textquotedblleft{}a shield\textquotedblright{} (CDIAL 5583). CDIAL marks the Hindi word masculine; standard Hindi today treats it as feminine.
\end{cousinweb}

\begin{grammarlens}[title={What you've built}]
One more word for telling what happened.
\end{grammarlens}

\subsection*{Your turn}
\begin{itemize}
\raggedright
  \item \emph{Say it:} \emph{ḍhāl}
  \item \emph{Say it:} \emph{ḍhāl}, once more
  \item \emph{Say it:} \emph{sipāhī ke ek hāth meṁ talvār aur dūsre meṁ ḍhāl thī.} — The soldier had a sword in one hand and a shield in the other.
  \item \emph{Recall:} say the Hindi for to fight hard, then the Hindi for a sword, then say \emph{ḍhāl} again
\end{itemize}

\begin{tcolorbox}[breakable,colback=teal!4,colframe=teal!35!black,title={Before you move on}]
What does \hi{ढाल} mean? (\textquotedblleft{}A shield\textquotedblright{}.) Say it once more.
\end{tcolorbox}
\section[\texorpdfstring{\hi{भाला} (bhālā)}{भाला (bhālā)}]{\hi{भाला} (bhālā) — a spear}

\label{lesson:HI-C322-bhala}



\begin{quote}
Before the new one: say the Hindi for a sword, then the Hindi for a shield.
\end{quote}

\subsection*{You'll want to know: \hi{भाला}}
\textbf{\hi{भाला}} — \emph{bhālā} — \textquotedblleft{}a spear\textquotedblright{}.

\begin{quote}
\hi{उस} \hi{सिपाही} \hi{के} \hi{हाथ} \hi{में} \hi{एक} \hi{लंबा} \hi{भाला} \hi{था।}
\end{quote}

\begin{quote}
\emph{us sipāhī ke hāth meṁ ek lambā bhālā thā.}
\end{quote}

\begin{quote}
That soldier had a long spear in his hand.
\end{quote}

Masculine.

\begin{cousinweb}
\textbf{\hi{भाला}} (\emph{bhālā}) goes back to Sanskrit \emph{bhalla} (CDIAL 9409), where CDIAL records Hindi \emph{bhālā}, a large spear.
\end{cousinweb}

\begin{grammarlens}[title={What you've built}]
One more word for telling what happened.
\end{grammarlens}

\subsection*{Your turn}
\begin{itemize}
\raggedright
  \item \emph{Say it:} \emph{bhālā}
  \item \emph{Say it:} \emph{bhālā}, once more
  \item \emph{Say it:} \emph{us sipāhī ke hāth meṁ ek lambā bhālā thā.} — That soldier had a long spear in his hand.
  \item \emph{Recall:} say the Hindi for a sword, then the Hindi for a shield, then say \emph{bhālā} again
\end{itemize}

\begin{tcolorbox}[breakable,colback=teal!4,colframe=teal!35!black,title={Before you move on}]
What does \hi{भाला} mean? (\textquotedblleft{}A spear\textquotedblright{}.) Say it once more.
\end{tcolorbox}
\section[\texorpdfstring{\hi{तीर} (tīr)}{तीर (tīr)}]{\hi{तीर} (tīr) — an arrow}

\label{lesson:HI-C322-tir}



\begin{quote}
Before the new one: say the Hindi for a shield, then the Hindi for a spear.
\end{quote}

\subsection*{You'll want to know: \hi{तीर}}
\textbf{\hi{तीर}} — \emph{tīr} — \textquotedblleft{}an arrow\textquotedblright{}.

\begin{quote}
\hi{उसने} \hi{एक} \hi{ही} \hi{तीर} \hi{से} \hi{फल} \hi{गिराया।}
\end{quote}

\begin{quote}
\emph{usne ek hī tīr se phal girāyā.}
\end{quote}

\begin{quote}
He brought down the fruit with a single arrow.
\end{quote}

Masculine.

\begin{cousinweb}
\textbf{\hi{तीर}} (\emph{tīr}) is usually traced to Persian \emph{tīr}, \textquotedblleft{}an arrow\textquotedblright{}. The bow is \textbf{\hi{धनुष}}; \textbf{\hi{तीर} \hi{चलाना}} is to shoot an arrow.
\end{cousinweb}

\begin{grammarlens}[title={What you've built}]
One more word for telling what happened.
\end{grammarlens}

\subsection*{Your turn}
\begin{itemize}
\raggedright
  \item \emph{Say it:} \emph{tīr}
  \item \emph{Say it:} \emph{tīr}, once more
  \item \emph{Say it:} \emph{usne ek hī tīr se phal girāyā.} — He brought down the fruit with a single arrow.
  \item \emph{Recall:} say the Hindi for a shield, then the Hindi for a spear, then say \emph{tīr} again
\end{itemize}

\begin{tcolorbox}[breakable,colback=teal!4,colframe=teal!35!black,title={Before you move on}]
What does \hi{तीर} mean? (\textquotedblleft{}An arrow\textquotedblright{}.) Say it once more.
\end{tcolorbox}
\section[\texorpdfstring{\hi{प्रायः} (prāyaḥ)}{प्रायः (prāyaḥ)}]{\hi{प्रायः} (prāyaḥ) — for the most part; usually}

\label{lesson:HI-C322-prayah}



\begin{quote}
Before the new one: say the Hindi for a spear, then the Hindi for an arrow.
\end{quote}

\subsection*{You'll want to know: \hi{प्रायः}}
\textbf{\hi{प्रायः}} — \emph{prāyaḥ} — \textquotedblleft{}for the most part; usually\textquotedblright{}.

\begin{quote}
\hi{उस} \hi{समय} \hi{युद्ध} \hi{प्रायः} \hi{तलवार} \hi{और} \hi{भाले} \hi{से} \hi{होते} \hi{थे।}
\end{quote}

\begin{quote}
\emph{us samay yuddh prāyaḥ talvār aur bhāle se hote the.}
\end{quote}

\begin{quote}
In those days wars were mostly fought with swords and spears.
\end{quote}

Formal and written; speech says \hi{अक्सर}.

\begin{cousinweb}
\textbf{\hi{प्रायः}} (\emph{prāyaḥ}) is Sanskrit \emph{prāyas}, \textquotedblleft{}for the most part, mostly\textquotedblright{} (Monier-Williams), and keeps the final sign \textbf{\hi{ः}}. It belongs to writing and to formal talk.
\end{cousinweb}

\begin{grammarlens}[title={What you've built}]
One more word for telling what happened.
\end{grammarlens}

\subsection*{Your turn}
\begin{itemize}
\raggedright
  \item \emph{Say it:} \emph{prāyaḥ}
  \item \emph{Say it:} \emph{prāyaḥ}, once more
  \item \emph{Say it:} \emph{us samay yuddh prāyaḥ talvār aur bhāle se hote the.} — In those days wars were mostly fought with swords and spears.
  \item \emph{Recall:} say the Hindi for a spear, then the Hindi for an arrow, then say \emph{prāyaḥ} again
\end{itemize}

\begin{tcolorbox}[breakable,colback=teal!4,colframe=teal!35!black,title={Before you move on}]
What does \hi{प्रायः} mean? (\textquotedblleft{}For the most part; usually\textquotedblright{}.) Say it once more.
\end{tcolorbox}
